\section{AlgoTune 探索：卷积与图计算的实现优化}\label{sec:mac}
\subsection{运行范围与结果}
Mac 探索使用固定 gpt-6.1-sol，在一维卷积和 PageRank 上各运行反馈与独立探索一次。每项研究最多 180 秒、5 次请求、3 个候选，使用两例开发输入；每例预热一次并测量五次。四轮都在提交第 3 个候选后结束，共 12 次模型调用、12 个生成候选，其中 11 个通过开发验证。

整批墙钟耗时 567.411 秒，四轮累计时长为 567.385 秒（其中模型等待 533.486 秒，开发评测 33.796 秒）。本批探索仅使用 2 例开发输入，配置中未启用最终测试。

\begin{smalltable}
\begin{tabularx}{\linewidth}{@{}Lrrrr@{}}\toprule
任务 / 策略 & 最佳候选 & 参考耗时（ms） & 候选耗时（ms） & 加速比\\\midrule
卷积 / F & c2 & 0.962906 & 0.741508 & 1.298578$\times$\\
卷积 / I & c1 & 1.117265 & 0.418563 & 2.669288$\times$\\
PageRank / F & c1 & 11.147868 & 0.445168 & 25.041935$\times$\\
PageRank / I & c1 & 10.908644 & 0.473231 & 23.051415$\times$\\\bottomrule
\end{tabularx}
\end{smalltable}

表中的函数时间先将两个用例的同次重复求和，再取五次重复的中位数。每次候选评估都会重新测量参考程序，因此不同候选的参考时间会随当时运行状态变化。

\reportfigure[0.88\linewidth]{mac_timings.pdf}{最佳候选及配对基线的计时。点为五次重复的中位数，区间为最小值至最大值；样本为两个用例的函数时间之和。卷积使用线性横轴，PageRank 使用对数横轴。}{fig:mac}

\subsection{卷积：执行路径、内存与精度}
卷积输入是长度 30000 和 8000 的两个 float64 数组，输出完整卷积的 37999 个元素。基线通过 \texttt{scipy.signal.convolve} 的通用入口完成计算；数值校验要求相对二范数误差满足
\[
 \frac{\|\hat y-y\|_2}{\|y\|_2+10^{-12}}\le 10^{-6}.
\]
这一容差为候选在数据表示、变换执行和缓冲区使用上提供了优化空间。

反馈组最佳候选 c2 直接使用实数快速傅里叶变换。它通过 \texttt{next\_fast\_len} 选择变换长度，将两个正向变换合并成双行批处理，并使用两个工作线程；逆变换阶段允许复用缓冲区。该方案针对的是通用接口之后的变换调度和内存路径，在本批输入上达到 1.298578 倍加速。

独立探索组的最佳 c1 在幅度与能量检查通过时使用 float32 实数 FFT，再将输出转换回 float64；特殊输入走双精度路径。降低中间表示精度减少了变换和内存开销，两个开发输入均通过给定容差，加速比为 2.669288 倍。

后续尝试提供了清晰的取舍。反馈组 c3 转向 overlap-add 分块卷积，输出正确，但加速比降到 0.699616，选择器继续保留 c2。独立组 c3 尝试导入该环境中不存在的底层 FFT 模块，记录为运行错误。开发评测既验证优化收益，也识别环境依赖和实现路径的代价。

\subsection{PageRank：用稀疏表示缩短执行路径}
两例 PageRank 输入各包含 1000 个节点，边数分别为 7950 和 7878。参考实现逐条构建 NetworkX 图对象，再调用 PageRank。候选将邻接表直接转换为 SciPy 稀疏矩阵，按出度归一化，以稀疏矩阵与向量乘法执行幂迭代，并显式处理无出边节点的概率质量。阻尼系数保持 0.85，最大迭代次数为 100，终止条件为一范数变化小于 $n\times10^{-6}$。

两个策略的最佳程序都缩短了“输入邻接表到迭代计算”之间的执行路径。独立组使用转置后的 CSR，反馈组以稀疏转置矩阵完成乘法；对应加速比为 23.051415 与 25.041935。这一结果的主要实现变化是直接使用适合计算的稀疏结构，减少图对象构建及格式转换。

反馈组后续 c2 进一步尝试把阻尼系数并入矩阵，并使用底层稀疏内核。其候选中位耗时为 0.391211 ms，低于 c1 的 0.445168 ms；但同次参考耗时也下降，因此比值为 24.919586，略低于 c1。由于每次评估均重测参考耗时，结合参考耗时、候选耗时与加速比共同观察，可以更准确地反映选择器的保留判定。

\subsection{候选轨迹与计时分离}
\begin{smalltable}
\begin{tabularx}{\linewidth}{@{}Lrrr@{}}\toprule
任务 / 策略 & c1 加速比 & c2 加速比 & c3 加速比\\\midrule
卷积 / F & 1.208234 & 1.298578 & 0.699616\\
卷积 / I & 2.669288 & 2.431732 & 运行错误\\
PageRank / F & 25.041935 & 24.919586 & 23.649574\\
PageRank / I & 23.051415 & 20.143776 & 21.508532\\\bottomrule
\end{tabularx}
\end{smalltable}

四个最佳候选的最大计时样本均小于对应参考程序的最小样本。例如，卷积反馈组的候选范围为 0.728046--0.781354 ms，参考范围为 0.907515--1.122496 ms；PageRank 反馈组分别为 0.413417--0.482418 ms 和 9.946698--11.152702 ms。两类任务均呈现出清晰的计时样本分离（完整范围见 \texttt{mac-analysis.json}）。

本批输入用量为 34240 tokens，输出为 26717 tokens；按冻结输入／输出价每百万 tokens 2／10 美元重算，费用为 USD 0.335650。四轮分别消耗 155.555、152.859、149.326、109.645 秒，生成和模型等待占据了主要研究历时。

\source{[S5] Mac 探索归档的 comparison、逐轮 report、ledger 与计时记录；归档校验说明见附录。}
